\section{Reproducibility Guide}
\label{appendix:reproducibility}

All code, trained models, datasets, and evaluation scripts are available at
\url{https://github.com/JDRanpariya/ball-balancing-on-arc}, together with
per-subdirectory READMEs that give exact run instructions. This appendix
only summarizes what is where; please refer to the repository for commands.

\subsection{Hardware}

Table~\ref{tab:hw_components} lists the physical components of the
ball-on-arc platform.

\begin{table}[H]
\centering
\caption{Hardware components.}
\label{tab:hw_components}
\begin{tabular}{lp{5.2cm}}
\toprule
\textbf{Component} & \textbf{Specification} \\
\midrule
3D-printed cart & PLA; its top surface is the convex arc ($R = 2.101$\,m, $\approx$30\,cm span) with a shallow central dip (2\,mm deep, 20\,mm FWHM); rides the actuator carriage \\
Ball & Steel, $m = 24$\,g, $r = 9$\,mm \\
Linear actuator & SMC LEFB32NXS-1000, 1000\,mm belt-driven stroke~\cite{smc_lefb_catalogue} \\
Servo motor \& drive & Beckhoff AM8121 motor via an EL7201 EtherCAT servo terminal \\
Controller & Beckhoff C6015 IPC, TwinCAT~3 runtime \\
Ball sensors & 2$\times$ VL53L0X ToF~\cite{vl53l0x_datasheet} at the arc edges (read alternately), plus a center IR beam-break, acquired by an Arduino Uno \\
Host PC & Windows (TwinCAT), Python stack in WSL2; USB serial to the drive and Arduino \\
\bottomrule
\end{tabular}
\end{table}

\paragraph{Note on length units.}
The rail provides approximately 1\,m of physical cart travel
(1000\,mm actuator stroke), yet the drive reports positions at roughly
twice that physical displacement, and the soft-limited operating window
spans 1.553\,m in these axis units ($\pm 0.7765$\,m about the center).

The factor comes from the drive's encoder scaling. The released NC
configuration (\texttt{hardware/twincat/axis\_parameters.xml}) sets
\texttt{ScaleFactorNumerator} $= 104.8576$ against
\texttt{ScaleFactorDenominator} $= 1048576$, that is $10^{-4}$ axis
units per encoder increment. With the EL7201 at its default 20
single-turn bits ($2^{20}$ increments per revolution), this amounts to
a feed constant of $104.8576$\,mm per motor revolution, which is the
encoder count rescaled rather than a measured quantity. The belt drive
advances $54$\,mm per revolution (the catalogued equivalent lead of the
LEFB32 AC-servo variant), so reported positions run
$104.8576/54 \approx 1.94$ times the physical ones. On that factor the
operating window is about $0.80$\,m of the $1000$\,mm stroke, leaving
roughly $0.10$\,m of margin at each end, which is what the mechanical
end stops require.

It affects
no result: all positions, velocities, limits, identified dynamics, and
results in the paper and code use these axis units consistently. The
factor matters only when comparing quoted track lengths against the
physical bill of materials.

\subsection{Software}

The reference environment pins Python~3.11, NumPy~2.2.6, SciPy~1.16.0,
PyTorch~\cite{paszke2019pytorch} 2.10.0, Stable-Baselines3~\cite{stable_baselines3} 2.9.0,
sb3-contrib~2.9.0, d3rlpy~\cite{seno2022d3rlpy} 2.8.1 (IQL), CasADi~\cite{casadi} 3.7.2 (NMPC),
and OSQP~\cite{osqp} 1.1.3 (MPC); exact versions are locked in \texttt{env.yaml} and
\texttt{requirements.lock.txt}.

\subsection{Repository layout}

\begin{itemize}
  \item \texttt{balancer/} (core library): all 13 controllers, simulation and hardware environments, and PLC/sensor interfaces.
  \item \texttt{hardware/}: bill of materials and assembly guide, 3D-print STLs, the TwinCAT project, Arduino sensor firmware, and wiring documentation.
  \item \texttt{training/}: RL, offline-RL, and world-model training scripts and configuration files.
  \item \texttt{evaluation/}: benchmark evaluation, controller tuning, and profiling scripts.
  \item \texttt{data/}: datasets and data-driven models (including the 1M random-exploration dataset and the ${>}$1M-transition logged-trial archive).
  \item \texttt{paper/}: paper source and the result data backing every table and figure.
  \item \texttt{scripts/}: repository utilities.
\end{itemize}

\subsection{Sensor calibration (hardware runs only)}

Reproducing figures and tables from the committed data needs no calibration.
For fresh on-hardware campaigns, the two time-of-flight sensors are
recalibrated first, following Stage~4 of the repository reproducibility guide:
the calibration script prints the per-sensor center offsets, which are
written back to \texttt{balancer/balancer/hardware/constants.py} once the two arc
halves report a consistent zero (the deployed datasets use the recalibrated
constants $189$/$158$~mm; the earlier $175$/$170$~mm values and the bias
they left are analyzed in
Appendix~\ref{appendix:sensor_characterization}). The offsets drift slowly, so we recalibrate before each new
data-collection or evaluation campaign.

\subsection{Reproducing the project}

The repository is built to reproduce the whole project end to end, not only
the paper's numbers. The step-by-step guide, \texttt{reproducibility\_guide.md},
with the per-directory READMEs, documents the full pipeline: building and
wiring the platform, deploying the PLC firmware, setting up the pinned
environment, calibrating the sensors, collecting datasets, training the
learning-based controllers, and evaluating all 13 controllers in simulation
and on the rig. The paper's tables and figures, the simulation validation,
and the wall-override ablations are the final stage of that pipeline and
regenerate from the committed data through the documented scripts, with no
hardware required. The master hardware-results file that backs Table~1 and the
hardware figures is itself rebuildable from the per-controller run aggregates
by \texttt{data/build\_consolidated.py}, which also emits the full
source-to-entry provenance manifest and checks that the rebuild is
byte-identical to the committed file. We do not duplicate the commands here.

\subsection{Companion website}

An interactive version of this benchmark, including the browser-based
simulation and per-controller trial videos, is available at
\url{https://ball-on-arc.pages.dev/}. The site carries no author,
affiliation, or institution identifiers.

\subsection{Key Parameters}

Table~\ref{tab:critical_params} collects the experimental parameters
held fixed across every benchmark trial.

\begin{table}[H]
\centering
\caption{Critical experimental parameters.}
\label{tab:critical_params}
\begin{tabular}{ll}
\toprule
\textbf{Parameter} & \textbf{Value} \\
\midrule
Control frequency & 20\,Hz (50\,ms period) \\
Trial duration & 30\,s maximum \\
Settling band & $\pm 0.01$\,rad \\
Settling duration & 1.0\,s continuous \\
Random seed (ICs) & 42 \\
Number of trials & 50 per controller \\
ToF timing budget & 34.5\,ms \\
Action range & $[-0.9, 0.9]$\,m/s (velocity set-point, axis units) \\
Cart limits & $\pm 0.7765$\,m (axis units) \\
\bottomrule
\end{tabular}
\end{table}
